\section{Formalization Target and Verified Results}\label{sec:formalization}

% Scope (PLAN.md §7.2): definition of weak solutions, the two capstone theorems,
% exact scope, and the main analytic difficulties.
% Statements MUST match claims/formalization-scope.md (pinned at
% leray-hopf v0.1.0-rc1, commit 7c15710a). No claim beyond that file.

The formalization target is the classical existence theorem for
Leray--Hopf weak solutions of the homogeneous, incompressible
Navier--Stokes equations~\cite{leray1934,hopf1951}: given a
divergence-free initial datum $u_0 \in L^2$, a viscosity $\nu > 0$, and a
time horizon $T > 0$, a weak solution on $[0,T]$ exists. We report a
Lean~4~\cite{moura2021lean4} formalization of this theorem, built on
mathlib~\cite{mathlib2020}, on two domains: the periodic unit 3-torus
$\mathbb{T}^3$ and whole space $\mathbb{R}^3$. All statements below are
quoted from the pinned reference commit of the public repository
\texttt{uda-lab/leray-hopf}, release \texttt{v0.1.0-rc1}, commit
\texttt{7c15710a7b9068a2aa105fc7c11b432e7685b7b5}~\cite{lerayhopf2026repo};
the repository's own scope statement (\texttt{README.md} at that commit)
is treated as the upper bound on what this section may claim.
% CLAIM: CLM-001 (leray-hopf@7c15710a7b9068a2aa105fc7c11b432e7685b7b5)

\subsection{Weak solutions as formalized}

A solution is a proof-carrying structure,
\texttt{LerayHopfSolutionFull} (on $\mathbb{T}^3$) and
\texttt{LerayHopfSolutionFull\_R3} (on $\mathbb{R}^3$), parameterizing a
curve $u : \mathrm{Time} \to D_\sigma$ together with the fields in
Table~\ref{tab:solution-fields}; both are definitional abbreviations of a
single domain-neutral \texttt{Galerkin.LerayHopfSolution} structure shared
by the two lanes. Two restrictions of the formalized notion are worth
flagging explicitly, since they are easy to over-read from the field
names alone. First, \texttt{weak\_eq} is tested only against
\emph{separated-variable} test functions $\psi(t)\!\cdot\! w(x)$, for a
scalar temporal factor $\psi$ and a fixed spatial test vector $w$ drawn
from the domain's divergence-free test class -- not a general space-time
test function $\varphi(t,x)$; the relation between the two formulations
is out of scope, neither proved nor assumed here. Second,
\texttt{initial\_trace} is a one-sided limit $u(t) \to u_0$ in the ambient
$L^2$ norm as $t \to 0^+$, not a weak-continuity statement on all of
$[0,T]$.

\begin{table}[h]
\centering
\caption{Solution-structure fields realizing the weak formulation.}
\label{tab:solution-fields}
\begin{tabular}{@{}p{0.20\linewidth}p{0.70\linewidth}@{}}
\toprule
Field & What is guaranteed \\
\midrule
\texttt{u} & valued in the closed divergence-free subspace $D_\sigma$ at every $t$ (type-level) \\
\texttt{weak\_eq} & weak Navier--Stokes identity against separated-variable tests $\psi(t)w(x)$, no forcing term \\
\texttt{energy\_ineq} & $\tfrac12\|u(t)\|^2 + \int_0^t \mathrm{dissip}(\nu,u(s))\,ds \le \tfrac12\|u_0\|^2$ for $t\in[0,T]$ \\
\texttt{initial\_trace} & $u(t)\to u_0$ in $L^2$ as $t\to 0^+$ (one-sided trace at $t=0$ only) \\
\texttt{energy\_class} & a.e.-in-$t$ membership in the $H^1$ regularity class, plus interval-integrable viscous dissipation \\
\texttt{u\_aestronglyMeasurable} & a.e.-strong measurability of $t \mapsto u(t)$ into the ambient $L^2$ space \\
\bottomrule
\end{tabular}
\end{table}

\subsection{Capstone existence theorems}

The two capstone statements are quoted verbatim from the pinned commit
(proof bodies elided):

\begin{lstlisting}
-- LerayHopf/Torus/GalerkinODECapstone.lean  (T^3)
theorem exists_lerayHopf_torus3 (u₀ : L2Sigma) (ν : ℝ) (hν : 0 < ν)
    (T : ℝ) (hT : 0 < T) :
    ∃ F : Torus3NSForms, Nonempty (LerayHopfSolutionFull F ν T u₀)

-- LerayHopf/R3/GalerkinODECapstone.lean     (R^3)
theorem exists_lerayHopf_r3 (u₀ : L2Sigma_R3) (ν : ℝ) (hν : 0 < ν)
    (T : ℝ) (hT : 0 < T) :
    ∃ (𝔊 : R3GalerkinScheme) (F : R3NSForms 𝔊),
    Nonempty (LerayHopfSolutionFull_R3 𝔊 F ν T u₀)
\end{lstlisting}
% CLAIM: CLM-001 (decl:exists_lerayHopf_torus3, decl:exists_lerayHopf_r3)

Both are unconditional in their stated hypotheses: for \emph{every}
divergence-free $u_0$, viscosity $\nu>0$, and horizon $T>0$, a solution
exists. On $\mathbb{R}^3$ the witness scheme $\mathfrak{G}$ (a concrete
Galerkin scheme built from a Schwartz divergence-free basis) is itself
produced by the proof, not assumed.

\subsection{Scope and non-claims}

The theorems above are exactly what is proved -- no more. This
formalization does \emph{not} claim: smoothness or any regularity beyond
the $H^1$ energy class exhibited by \texttt{energy\_class}; uniqueness or
non-uniqueness of the constructed solution; an equation with an external
forcing term (the homogeneous case only); a single solution valid
uniformly on $[0,\infty)$ (each capstone fixes a finite horizon $T$, taken
as an arbitrary but fixed input); or the weak formulation against general
space-time test functions (only the separated-variable class
$\psi(t)\!\cdot\! w(x)$, as above). A further, purely engineering scope
boundary: the repository contains an explicit
\texttt{LerayHopf.Experimental} namespace holding incomplete
Bochner-space-in-time developments that neither capstone depends on and
that are excluded from the release surface described in
Section~\ref{sec:release-surface} below.

\subsection{Main analytic difficulties}

The 96-file, $\sim$42{,}000-line \texttt{LerayHopf/} tree (measured at the
pinned commit; \texttt{evidence/metrics/formalization-metrics.json})
% CLAIM: CLM-001 (formalization-metrics.json, measured at leray-hopf@7c15710a)
is organized around three domain-neutral layers shared by both lanes, plus
per-domain instantiations. \emph{Galerkin approximation and finite-dimensional ODE
existence}: for each finite-dimensional Galerkin subspace $V_n$, the
truncated quadratic vector field is shown $C^1$ with a dissipative
inner-product estimate, and forward-global existence on $V_n$ is obtained
by a tiling/gluing argument over local Picard--Lindelöf solutions driven
by an a~priori energy bound, rather than by invoking a global
existence theorem from mathlib directly. \emph{Spatial compactness}:
$\mathbb{T}^3$ uses a Fourier mode-tail Rellich argument; $\mathbb{R}^3$
uses a local Rellich--Kondrachov embedding on balls (the one classical
input not otherwise derived from mathlib) combined with a
Fréchet--Kolmogorov criterion and a countable diagonal argument over an
exhausting sequence of balls. \emph{Aubin--Lions compactness in time and
limit passage}: a shared Gelfand-triple/Bochner-space layer supplies the
time-compactness and weak-limit toolkit consumed by both lanes' passage
from the Galerkin sequence to a solution of the weak equation, energy
inequality, and initial trace; the $\mathbb{T}^3$ route additionally
exploits a mode-wise spectral decomposition, while the $\mathbb{R}^3$
route uses Steklov time-averaging. The nonlinear convection term on each
domain is handled by a proof-carrying trilinear extension pinned to the
Schwartz/Galerkin test class, not by a claim of a bounded convection
operator on all of $L^2\times L^2\times L^2$.

\subsection{Kernel-only axiom check}

Both capstones are \emph{kernel-only}: \texttt{\#print axioms} on
\texttt{exists\_lerayHopf\_torus3} and \texttt{exists\_lerayHopf\_r3}
returns exactly the three standard Lean kernel axioms --
\texttt{propext}, \texttt{Classical.choice}, \texttt{Quot.sound} -- with
zero project-declared axioms and no \texttt{sorryAx} on either theorem.
This is pinned by a CI script (\texttt{scripts/check-axioms-live.sh})
that runs the live \texttt{\#print axioms} check and fails on any
unexpected or missing axiom.

\subsection{Release-surface enforcement and build attestation}\label{sec:release-surface}

\texttt{import LerayHopf} is defined as the complete release surface:
both capstones plus every supporting file they depend on. A dedicated
guardrail script, \texttt{scripts/check-release-cone.sh}, walks the
transitive import closure of that entry point and fails if any reachable
file contains a \texttt{sorry}, a project \texttt{axiom}/\texttt{opaque},
or a placeholder namespace; this check runs as a fast, buildless textual
guardrail on every pull request. The \texttt{LerayHopf.Experimental}
namespace mentioned above is, by this same enforcement, unreachable from
the release surface. A separate, heavier step -- a full \texttt{lake
build} together with the live axiom pin -- is run manually, on demand,
against a specific commit, and recorded as a release-candidate build
attestation; it certifies exactly the attested commit SHA, not
continuously the branch tip. The results reported in this section are
attested against commit \texttt{7c15710a} under release tag
\texttt{v0.1.0-rc1} in exactly this manner, not inferred from a
continuously green push-triggered build.
